\section{Jump Game}
\label{sec:jump-game}

\problemstatement{We are given an array $\textit{nums}$ of length $n$, where
$\textit{nums}[i]$ is the maximum number of steps we are allowed to jump
forward from index $i$ (we may jump anywhere from $1$ up to
$\textit{nums}[i]$ steps, or choose to jump fewer). Starting at index $0$, we
must decide whether it is possible to reach index $n-1$ \textbf{at all},
following any sequence of valid jumps. Return \texttt{true} or
\texttt{false}.}

\begin{patternbox}
\emph{``Can you reach the last index, given a maximum jump length at each
position''} is the defining example of a \emph{reachability-horizon}
greedy. The keyword to notice is that we only need a yes/no feasibility
answer, not the actual sequence of jumps or the minimum number of jumps --
that distinction matters, because it means we never need to track \emph{how}
we got somewhere, only the single number \emph{how far forward can we
possibly see from everything visited so far}. Whenever a problem asks
``can position $X$ be reached'' under a rule where each position extends
your reach by some amount, track a running \emph{farthest reachable index}
and sweep once, left to right.
\end{patternbox}

\subsection*{Understanding the problem carefully}

Define $\textit{farthest}$ as the largest index reachable using any
combination of jumps from indices we have already confirmed are reachable,
considering only indices up to the one we are currently examining. As we
scan left to right through index $i = 0, 1, \dots, n-1$:

\begin{itemize}
  \item If $i \le \textit{farthest}$, then index $i$ itself is reachable
        (it lies within the horizon established by earlier indices), so we
        are allowed to jump from it: $\textit{farthest}$ can potentially be
        extended to $i + \textit{nums}[i]$.
  \item If $i > \textit{farthest}$, then index $i$ is \emph{not} reachable
        by anything before it, and since we are scanning strictly left to
        right, nothing later can make index $i$ reachable either -- the
        whole array up to this point has failed to bridge the gap, so the
        answer is immediately \texttt{false}.
\end{itemize}

\begin{ideabox}[Greedy Choice]
Maintain a single running value $\textit{farthest}$, the best (largest)
index reachable using everything examined so far. At each index $i$, first
check whether $i$ itself is still within $\textit{farthest}$ (if not, fail
immediately); otherwise, update
$\textit{farthest} \gets \max(\textit{farthest},\, i + \textit{nums}[i])$.
If at any point $\textit{farthest} \ge n-1$, the last index is reachable and
we may stop early and answer \texttt{true}.
\end{ideabox}

\subsection*{Why tracking only the farthest reach is correct}

This works because reachability is monotonic in a useful sense: if index
$j \le \textit{farthest}$ is reachable, every index from $0$ up to
$\textit{farthest}$ is also reachable, because we can always choose to jump
\emph{short} of our maximum allowance at any earlier reachable index to land
exactly where we want (jumping any number of steps from $1$ up to the
maximum is allowed, so every intermediate landing spot is achievable, not
just the farthest one). This means we never need to track the \emph{set} of
reachable indices, nor which specific path of jumps was used -- only the
single supremum of that set, $\textit{farthest}$, fully describes everything
we need to know about reachability at each step. There is no benefit to
"saving" a long jump for later or being clever about which specific jump
length to use, because every index up to the maximum reach is already
reachable regardless of which path got us there. This collapses what could
look like an exponential search over jump sequences into one linear pass
that only ever needs to remember one number.

\subsection*{Algorithm}

\begin{enumerate}
  \item Initialize $\textit{farthest} \gets 0$.
  \item For $i \gets 0$ to $n-1$:
  \begin{enumerate}
    \item If $i > \textit{farthest}$: return \texttt{false} (index $i$, and
          therefore the rest of the array, is unreachable).
    \item $\textit{farthest} \gets \max(\textit{farthest},\, i +
          \textit{nums}[i])$.
    \item If $\textit{farthest} \ge n-1$: return \texttt{true} early.
  \end{enumerate}
  \item If the loop completes without returning \texttt{false}, return
        \texttt{true}.
\end{enumerate}

\begin{algorithm}[H]
\caption{Jump Game (Reachability)}
\begin{algorithmic}[1]
\Procedure{CanJump}{\textit{nums}}
  \State $\textit{farthest} \gets 0$
  \State $n \gets \text{length}(\textit{nums})$
  \For{$i \gets 0$ to $n-1$}
    \If{$i > \textit{farthest}$}
      \State \Return \textbf{false}
    \EndIf
    \State $\textit{farthest} \gets \max(\textit{farthest}, i + \textit{nums}[i])$
    \If{$\textit{farthest} \ge n - 1$}
      \State \Return \textbf{true}
    \EndIf
  \EndFor
  \State \Return \textbf{true}
\EndProcedure
\end{algorithmic}
\end{algorithm}

\subsection*{Dry run}

\dryrun{Take $\textit{nums} = [2, 3, 1, 1, 4]$ ($n=5$, last index $=4$).}

\begin{itemize}
  \item $i=0$: $0 \le \textit{farthest}=0$. Update
        $\textit{farthest} = \max(0, 0+2) = 2$.
  \item $i=1$: $1 \le 2$. Update
        $\textit{farthest} = \max(2, 1+3) = 4$. Since $4 \ge n-1 = 4$,
        return \texttt{true} immediately.
\end{itemize}

The answer is \texttt{true}: from index $0$ we can jump $2$ steps to index
$2$... but more directly, from index $1$ (reachable, since $1 \le 2$) a jump
of $3$ steps reaches index $4$, the last index.

Now take $\textit{nums} = [3, 2, 1, 0, 4]$ ($n=5$, last index $=4$).

\begin{itemize}
  \item $i=0$: $0 \le 0$. $\textit{farthest} = \max(0, 0+3) = 3$.
  \item $i=1$: $1 \le 3$. $\textit{farthest} = \max(3, 1+2) = 3$.
  \item $i=2$: $2 \le 3$. $\textit{farthest} = \max(3, 2+1) = 3$.
  \item $i=3$: $3 \le 3$. $\textit{farthest} = \max(3, 3+0) = 3$.
  \item $i=4$: $4 > \textit{farthest}=3$. Return \texttt{false}.
\end{itemize}

The answer is \texttt{false}: index $3$ has a jump length of $0$, and
everything before it can reach at most index $3$, so index $4$ -- and the
useful $4$ sitting right there at index $4$ -- can never be reached.

\subsection*{C++ implementation}

\begin{lstlisting}
#include <bits/stdc++.h>
using namespace std;

bool canJump(vector<int>& nums) {
    int n = nums.size();
    int farthest = 0;

    for (int i = 0; i < n; i++) {
        if (i > farthest) return false;
        farthest = max(farthest, i + nums[i]);
        if (farthest >= n - 1) return true;
    }
    return true;
}

int main() {
    vector<int> nums = {2, 3, 1, 1, 4};
    cout << (canJump(nums) ? "true" : "false") << endl;
    return 0;
}
\end{lstlisting}

\begin{complexitybox}
\textbf{Time Complexity.} We scan the array once, doing constant work per
index, so the time complexity is
\[
  O(n).
\]
\textbf{Space Complexity.} We maintain only the single variable
\textit{farthest}, so the space complexity is $O(1)$.
\end{complexitybox}

\begin{takeawaybox}
When a problem only asks \emph{whether} some destination is reachable under
a rule that extends your reach from any currently-reachable position, you
rarely need to enumerate paths. Track the single furthest horizon reachable
so far and sweep once -- this is strictly more efficient than recursion or
BFS/DP over individual positions, and it generalizes directly to the
minimum-jumps variant in the next section.
\end{takeawaybox}
